\section{Module \texttt{input}}

Le module \texttt{input} encapsule les widgets de saisie utilisateur: champ texte, liste déroulante et spin button.

\subsection{Macros de Configuration}

\begin{docbox}{Macros \texttt{ENTRY\_CONFIG\_DEFAULT}, \texttt{DROPDOWN\_CONFIG\_DEFAULT}, \texttt{SPIN\_BUTTON\_CONFIG\_DEFAULT}}
\begin{lstlisting}[language=C]
#define ENTRY_CONFIG_DEFAULT ((entry_config){ \
    .placeholder = NULL, \
    .default_text = NULL, \
    .purpose_hint = NULL, \
    .password = false, \
    .editable = true, \
    .max_length = 0, \
    .style = WIDGET_STYLE_CONFIG_DEFAULT, \
    .on_changed = NULL, \
    .user_data = NULL, \
})

#define DROPDOWN_CONFIG_DEFAULT ((dropdown_config){ \
    .options = NULL, \
    .selected_index = 0, \
    .enable_search = false, \
    .style = WIDGET_STYLE_CONFIG_DEFAULT, \
    .on_selected = NULL, \
    .user_data = NULL, \
})

#define SPIN_BUTTON_CONFIG_DEFAULT ((spin_button_config){ \
    .min_value = 0.0, \
    .max_value = 100.0, \
    .step = 1.0, \
    .page = 5.0, \
    .value = 0.0, \
    .digits = 0, \
    .numeric_only = false, \
    .wrap = false, \
    .style = WIDGET_STYLE_CONFIG_DEFAULT, \
    .on_value_changed = NULL, \
    .user_data = NULL, \
})
\end{lstlisting}
\end{docbox}

Ces macros fixent une base stable pour les champs de saisie et évitent les paramètres incohérents. Les callbacks nuls et pointeurs \texttt{NULL} empêchent des appels invalides au runtime. C'est une protection simple mais efficace contre les segfaults en C.

L'initialisation par macros rend aussi le comportement de saisie reproductible entre écrans: même valeur par défaut, même politique de validation, même branchement des événements. Cela simplifie les tests et la maintenance.

\subsection{Fonctions API}

\begin{lstlisting}[language=C]
GtkWidget *create_entry(const entry_config *config) {
    GtkWidget *entry = gtk_entry_new();
    if (config == NULL) return entry;

    if (config->placeholder) {
        gtk_entry_set_placeholder_text(GTK_ENTRY(entry), config->placeholder);
    }
    if (config->default_text) {
        gtk_editable_set_text(GTK_EDITABLE(entry), config->default_text);
    }

    gtk_entry_set_visibility(GTK_ENTRY(entry), !config->password);
    gtk_editable_set_editable(GTK_EDITABLE(entry), config->editable);
    apply_widget_style(entry, &config->style);

    if (config->on_changed) {
        g_signal_connect(entry, "changed", G_CALLBACK(config->on_changed), config->user_data);
    }
    return entry;
}
\end{lstlisting}

La logique de création sépare clairement la configuration visuelle, le comportement de saisie et la connexion des événements utilisateur.

Le module \texttt{input} formalise ainsi un flux de données propre: saisie utilisateur, transformation, puis notification via signal. Cette séparation est utile pour concevoir une interface robuste et évolutive.

\newpage
